% Gaya run journals and 2D reduced-model comparison, November 2026.
\paragraph{Full-order reference runs.}
The native Cartesian DG-P0 implementation shares cell and face evaluators
across one, two and three dimensions. Tests completed in each dimension on
one and six MPI ranks. The $250^3$ 3D case, with three velocity components,
advanced $46\,875\,000$ degrees of freedom through 500 time steps of size
$0.05$. The run used 256 MPI ranks on two Gaya CPU nodes and took
$5\,552$~s of elapsed time. Six field snapshots were exported.

\paragraph{Reduced-order comparison in 2D.}
The $250\times250$ DG-P0 case has $125\,000$ degrees of freedom. It uses a
Cartesian Godunov flux and trapezoidal time integration. The reduced models
share a 100-mode POD basis and were compared at the held-out parameter
$\mu=(4.75,0.02)$ over the same 500 time steps. Errors are relative $L^2$
errors over the full trajectory and both velocity components.

\begin{table}[!htb]
\centering
\small
\caption{Burgers 2D comparison on one Gaya CPU core. Timings cover the
online nonlinear solves.}
\label{tab:app-feelpp-nlcrb:gaya-2d}
\begin{tabular}{@{}lrrr@{}}
\toprule
\textbf{Method} & \textbf{Integration cells} & \textbf{Time (s)} &
\textbf{Trajectory error (\%)}\\
\midrule
Full-order reference & $62\,500$ & 256.18 & Reference\\
POD-Galerkin & $62\,500$ & 477.63 & 0.984\\
POD-Galerkin with ECSW & 300 & 2.84 & 1.008\\
POD-LSPG with ECSW & 298 & 6.39 & 0.984\\
\bottomrule
\end{tabular}
\end{table}

ECSW assigns positive weights to the selected cells and retains adjacent
cells needed to evaluate face terms. The Galerkin and LSPG samples use
closures of $1\,438$ and $1\,427$ cells, respectively. The measured
comparison establishes projection and hyper-reduction baselines for
the nonlinear compressive models.
